\clearpage
\fancyhead[L]{\small BESTIMMUNG DER ABSOLUTEN DECLINATION.}
\SourceLeaf{442}\pdfbookmark[0]{Über die Anwendung des Magnetometers zur Bestimmung der absoluten Declination.}{absolute-declination-title}
\begin{center}
{\Large\bfseries ÜBER DIE ANWENDUNG DES MAGNETOMETERS\par\medskip ZUR BESTIMMUNG DER ABSOLUTEN DECLINATION.}
\par\bigskip
Resultate aus den Beobachtungen des magnetischen Vereins. 1841. I.
\end{center}
Es ist nicht meine Absicht, den in der Überschrift bezeichneten Gegenstand, über welchen bereits im 2. Bande der \textit{Resultate} [1837. VII.] ein sehr ausführlicher Aufsatz mitgetheilt ist, hier noch einmal vollständig abzuhandeln. Ich werde vielmehr mich hier auf Eine Hauptaufgabe beschränken, in Beziehung auf welche die a. a. O. S. 121—124 gegebene Entwickelung als ungenügend erscheint: diese Aufgabe betrifft die Bestimmung des Azimuths derjenigen Verticalebene, in welcher sich die optische Achse des Beobachtungsfernrohrs befindet.

Die in Rede stehende Verticalebene ist festgestellt durch die Marke und einen festen Punkt der Scale, welcher durch den über der Mitte des Objectivs des Beobachtungsfernrohrs herabhangenden Lothfaden bestimmt wird. Von dem Standpunkte des Beobachtungsfernrohrs aus muss ein entfernter Gegenstand sichtbar sein, dessen Azimuth anderweitig schon bekannt ist, und es kommt also zunächst darauf an, den auf den Horizont projicirten Winkel zwischen diesem Gegenstande und der Marke zu bestimmen. Ich nehme an, dass zu diesem Geschäft ein Theodolith nach der bekannten von \textsc{Reichenbach} eingeführten Construction angewandt wird, ohne darum zugleich vorauszusetzen, dass derselbe Theodolith auch zu den magnetischen Beobachtungen gebraucht werde, wozu vielmehr füglich ein besonderes Ablesungsfernrohr verwandt werden kann.

Der gewöhnliche Gebrauch solcher Theodolithen bezieht sich auf Winkel\-messungen zwischen Gegenständen in so grossen Entfernungen, dass eine geringe
\expandafter\def\csname HistoricalNoteText443\endcsname{Bei den weiter unten anzuführenden Beobachtungen etwa 20 Millimeter.}%
\SourceLeaf{443}Abweichung von mehrern der Idee des Instruments zum Grunde liegenden Bedingungen in der Ausführung seines Baues einen merklichen Fehler nicht hervorbringen kann, wie denn in der That absolute Vollkommenheit in keiner mechanischen Arbeit erreichbar ist. Allein wenn die Gegenstände (oder wie im vorliegenden Falle einer derselben) vergleichungsweise sehr nahe sind, so wird es allerdings nothwendig, es mit solchen Abweichungen schärfer zu nehmen, und namentlich müssen hier folgende Umstände in Erwägung kommen.

I. Die verticale Drehungsachse, die horizontale Drehungsachse und die optische Achse des Fernrohrs sollten einander in Einem Punkte schneiden. In so fern dieser Bedingung vollkommen nicht genügt ist, wird eine dreifache Abweichung vorkommen. Es seien $A,B$ resp. die beiden Punkte in der verticalen und der horizontalen Drehungsachse, wo diese einander am nächsten sind; imgleichen $C,D$ die ähnlichen Punkte der horizontalen Drehungsachse und der optischen Achse. Man bezeichne die Entfernungen $AB,BC,CD$ mit $\alpha,\beta,\gamma$, unter beliebiger Bestimmung rücksichtlich der Zeichen; man mag z. B. $\alpha$ positiv setzen, wenn $A$ auf derselben Seite der horizontalen Drehungsachse liegt wie das Ocular des Fernrohrs; $\beta$ positiv, wenn für den am Ocular stehenden Beobachter der Punkt $C$ rechts von $B$ fällt; $\gamma$ positiv, wenn $D$ oberhalb $C$ fällt.

II. Die optische Achse des Fernrohrs sollte normal gegen die horizontale Drehungsachse sein. Dieser Bedingung kann man zwar mit aller nöthigen Schärfe Genüge leisten: allein da man, um nach der Beobachtung eines entfernten Gegenstandes einen nahen deutlich sehen zu können, nothwendig die Ocularröhre weiter\HistoricalNote{443} herausziehen, also dem Fadenkreuze eine veränderte Stellung gegen das Objectiv geben muss, so ist man nicht berechtigt vorauszusetzen, dass beiden Stellungen der Ocularröhre einerlei optische Achse entspreche, sondern muss darauf gefasst sein, dass die für eine Stellung gemachte Berichtigung bei der andern wieder verloren gehe. Grösserer Allgemeinheit wegen mag man voraussetzen, dass für keine von beiden Stellungen die Berichtigung genau gemacht sei, und den Collimationsfehler für die erste Stellung mit $c$, für die zweite mit $c'$ bezeichnen: als positiv mag man dieselben annehmen, wenn die optische Achse mit dem dem Beobachter rechts liegenden Arme der horizontalen Achse einen spitzen Winkel macht.

\SourceLeaf{444}Offenbar werden auch, wenn die Grössen $\beta,\gamma$ der erstern Ocularstellung angehören, etwas veränderte Werthe bei der zweiten an ihre Stelle treten, die mit $\beta',\gamma'$ bezeichnet werden mögen.

Es ist nun zwar leicht, den Einfluss aller dieser Abweichungen auf die Messung sowohl des horizontalen Winkels zwischen den beiden Gegenständen, als ihrer Elevationen (wenn der Theodolith zugleich einen Höhenkreis hat) in strengen Gleichungen darzustellen, aus welchen die Resultate vermittelst einer biquadratischen Gleichung abzuleiten sein würden; allein da die sieben Grössen $\alpha,\beta,\gamma,\beta',\gamma',c,c'$ alle nur sehr klein sein können, so kann man unbedenklich alle Grössen. welche in Beziehung auf jene von der zweiten oder höherer Ordnung sind, vernachlässigen, und das Resultat ihres Einflusses in sehr einfache Form bringen. Aber selbst dieser Darstellung können wir hier überhoben sein. Man sieht nemlich leicht ein, dass, wenn man das Fernrohr auf gehörige Art umlegt, sämmtliche sieben Abweichungen, ohne ihre Grösse zu ändern, bloss die entgegengesetzten Zeichen annehmen, und dass mithin dasselbe auch von den Fehlern der Messungen gelten wird, die man bei den zwei verschiedenen Arten des Einliegens anstellt. Das Mittel aus diesen beiden Messungen ist folglich von dem Einflusse dieser Fehler, ohne dass man die einzelnen Bestandtheile davon zu kennen braucht, von selbst befreit, und man erhält dadurch den wahren Werth des Winkels zwischen den beiden in der Verticalachse des Instruments sich schneidenden Verticalebenen, in denen die beiden Gegenstände liegen. Dasselbe gilt von den Elevationen, welche sich dann auf den Punkt $A$ beziehen, aber für unsern gegenwärtigen Zweck unnöthig sind.

Das Umlegen muss so geschehen, dass die Zapfen wieder in dieselben Pfannen zu liegen kommen, während die obere Seite des Fernrohrs zur untern wird und das Objectiv an die Stelle des Oculars kommt: es ist also dies Umlegen dasselbe, was eine halbe Umdrehung um die horizontale Achse sein würde, welche auszuführen die Stützen nur nicht hoch genug sind. Wollte man anstatt dieser Art das Umlegen so verrichten, dass die Zapfen in die andern Pfannen gelegt würden, während das Objectivende auf derselben Seite bliebe (was geometrisch betrachtet einerlei ist mit einer halben Umdrehung um die Achse des Fernrohrs), so würden nicht alle sieben Grössen $\alpha,\beta,\gamma,\beta',\gamma',c,c'$ in dem Fall sein, schlechthin die entgegengesetzten Zeichen anzunehmen, sondern dies würde nur von $\gamma,\gamma',c,c'$ gelten. Man hat nemlich keine Sicherheit, dass die Stützen \textit{genau}
\expandafter\def\csname HistoricalNoteText445\endcsname{Auf der ersten Tafel des ersten Bandes der \textit{Resultate} [1836. I.] ist dieser Thurm angedeutet, ungefähr so, wie er bei nicht geöffnetem Fenster von dem Theodolithenplatz aus erscheint: an dem Orte des Auges, welcher der perspectivischen Zeichnung eigentlich zum Grunde liegt, wird der Thurm durch die Wand links vom Fenster verdeckt.}%
\SourceLeaf{445}gleich weit von der Verticalachse abstehen, und es würden daher. nach solchem Umlegen, der Punkt $B$ ein anderer sein können als vorher, mithin auch $\beta$ und $\beta'$ andere Werthe annehmen. Dass zugleich $\alpha$ das entgegengesetzte Zeichen nicht annimmt, sondern ganz den vorigen Werth behält, ist übrigens allerdings hier unwesentlich, weil in dem linearen Ausdruck für den Fehler der horizontalen Winkelmessung $\alpha$ gar nicht vorkommt.

Wie nun eine solche Winkelmessung für den beabsichtigten Zweck zu benutzen sei, wird sich am einfachsten durch ein Beispiel zeigen lassen, wozu ich die letzte am 11. März d. J. ausgeführte Anwendung des Verfahrens wähle.

In dem hiesigen magnetischen Observatorium dient zur Anknüpfung der Beobachtungen an den wahren astronomischen Meridian ein Stadtkirchthurm, dessen Knopfstange an dem Platze des Beobachtungsfernrohrs durch das geöffnete nordliche Fenster frei sichtbar ist\HistoricalNote{445}, und zwar von der Mitte der Säule aus, welche seit Julius 1837 an die Stelle des früher gebrauchten hölzernen Stativs getreten ist, in dem Azimuth $173^0 35'25''5$. Gefunden war dieses Azimuth, indem man einen Theodolithen an einer andern Stelle des Saales aufstellte, die Verticalachse genau im Allignement der Mitte der Säule und des Kirchthurms, und die Winkel zwischen letzterm und zweien andern daselbst sichtbaren Kirchthürmen maass; die Lage dieser verschiedenen Thürme gegen den Nullpunkt in der Sternwarte war durch frühere an die Gradmessung geknüpfte Messungen genau bekannt, und das in Rede stehende Azimuth liess sich daher aus jenen Winkelmessungen leicht berechnen.

Es wurde nun ein achtzolliger \textsc{Ertel}scher Repetitionstheodolith auf der Säule so aufgestellt, dass seine Verticalachse so genau wie möglich mit der Mitte der Säule zusammenfiel, und der horizontale Winkel zwischen der Marke und der Knopfstange des Thurms bei den beiden verschiedenen oben bezeichneten Arten des Einliegens des Fernrohrs, jedesmal durch 25 Repetitionen, gemessen. In der ersten Lage fand sich der Winkel
\[
=11^0 40'54''50
\]
\SourceLeaf{446}in der zweiten
\[
=11^0 41'36''18
\]

Der wahre Werth des Winkels, seinen Scheitel in die Verticalachse des Theodolithen gesetzt, ist folglich
\[
=11^0 41'15''34
\]
mithin das Azimuth der durch diese Verticalachse und die Marke gelegten Verticalebene
\[
=161^0 54'10''16
\]

Mit dieser Operation war eine andere verbunden, deren Zweck war, auszumitteln, in welchem Punkte die Scale von dieser Verticalebene geschnitten wird.

Auf dem Objectivende des Theodolithenfernrohrs ist ein Ring aufgesteckt, der auf seiner Vorderfläche zwei einander diametral gegenüber liegende zarte Einschnitte und diesen correspondirend auf der äussern runden Fläche zwei Häkchen hat, in welche nach der verschiedenen Lage des Fernrohrs ein feiner mit einem Gewichte beschwerter Goldfaden eingehängt wird. Der Ring wird so gedreht, dass der durch die Einschnitte gehende Diameter gegen die horizontale Drehungsachse des Fernrohrs normal ist: man erkennt die Erfüllung dieser Bedingung, wenn der in dem obern Einschnitte einliegende Lothfaden zugleich genau dem untern entspricht, zu welchem Ende man das Fernrohr nahe horizontal stellen muss, nemlich nur so wenig nach unten geneigt, dass der Faden noch eben frei vor dem Ringe spielen kann: die Coincidenz wird mit einer Loupe geprüft. Der Lothfaden spielt in einer sehr geringen Entfernung vor der Scale, und es kommt nun darauf an, die correspondirenden Punkte der Scale in den beiden verschiedenen Lagen des Fernrohrs, indem es jedesmal auf die Marke, oder vielmehr in deren Verticalebene gerichtet ist, zu notiren. Genau genommen, sind damit diejenigen Punkte der Scale gemeint, welche in der durch die Marke und den Lothfaden gehenden Verticalebene liegen, und man kann dies unmittelbar in dem Spiegel des Magnetometers erkennen, wenn der Theodolith selbst die Bestimmung hat, als Ablesungsfernrohr zu dienen, also die Scale sich in einer dieser Bestimmung angemessenen Höhe befindet. Es ist wohl überflüssig zu erinnern, dass es in diesem Falle nothwendig werden kann, den Magnetstab
\SourceLeaf{447}des Magnetometers vermittelst eines aus der Ferne wirkenden Ablenkungsstabes erst in eine solche Stellung zu bringen, dass der betreffende Scalenpunkt nahe am Fadenkreuz des Theodolithenfernrohrs erscheint. Im hiesigen magnetischen Observatorium, wo jetzt die magnetischen Beobachtungen mit einem besondern Ablesungsfernrohre angestellt werden, welches sich in einer geringern Höhe über der Säule befindet als das Theodolithenfernrohr, ist mit diesem das Bild der in einer der Lage des Ablesungsfernrohrs angemessenen Höhe angebrachten Scale im Spiegel des Magnetometers nicht sichtbar. Ich habe daher zur Bestimmung des dem vom Theodolithenfernrohr herabgehenden Lothfaden correspondirenden Scalenpunktes das Ablesungsfernrohr selbst gebraucht, welches zu diesem Zweck nahe an der Marke in der betreffenden Verticalebene aufgestellt wurde: dass man nicht nöthig hat, wegen letzterer Bedingung gar zu ängstlich zu sein, in sofern der Lothfaden, wie schon bemerkt ist, in geringer Entfernung von der Scale spielt, leuchtet von selbst ein. Es fand sich auf diese Weise der Lothfaden correspondirend
\begin{quote}
dem Scalenpunkte $850.8$ bei der ersten Lage des Theodolithenfernrohrs, und

 dem Punkte $849.4$ bei der zweiten Lage,
\end{quote}
woraus man schliessen darf, dass die durch die Marke und die Verticalachse des Theodolithen gehende Verticalebene, deren Azimuth oben bestimmt ist, die Scale in dem Punkte $850.1$ schneidet.

Die Bestimmung des Azimuths derjenigen Verticalebene, in welcher sich die optische Achse des Beobachtungsfernrohrs befindet, hat nun weiter keine Schwierigkeit. Correspondirt der vor der Mitte des Objectivs desselben herabhängende Lothfaden dem Scalenpunkte $850{,}1+n$, so reicht es hin (weil die Scale als normal gegen jene Ebene gestellt vorausgesetzt wird), das Product $n.206265''$ mit der horizontalen Entfernung der Scale von der Marke, in Scalentheilen ausgedrückt, zu dividiren, und den Quotienten mit seinem Zeichen zu $161^0 54'10''16$ hinzuzufügen. Gegenwärtig ist jene Entfernung $=9638.7$. Diente also der Theodolith selbst, und zwar bei der ersten Lage des Fernrohrs, zum Beobachten, so wäre dieses Azimuth
\[
=161^0 54'25''1
\]
bei der zweiten Lage hingegen
\[
=161^0 53'55''2
\]
\SourceLeaf{448}Da aber, wie schon bemerkt ist, zum Beobachten ein besonderes Ablesungsfernrohr dient, welches nach der Beendigung der obigen Operationen so aufgestellt wurde, dass, bei der Richtung der optischen Achse auf die Verticale der Marke, der Lothfaden dem Punkte $850.0$ entsprach, so ist das verlangte Azimuth
\[
=161^0 54'8''0
\]

Es mögen über das hier behandelte Geschäft noch ein Paar Bemerkungen hier beigefügt werden.

I. Wenn die horizontale Achse in ihren Lagern einigen Spielraum in dem Sinn ihrer Länge hat, so muss man Sorge tragen, dass sie bei den einzelnen Winkelmessungen immer gleiche Lage gegen die Stützen habe, etwa dadurch, dass man jedesmal den Spielraum auf Einer Seite durch einen leichten Druck gegen das Ende eines bestimmten Zapfens zum Verschwinden bringt. Ohne diese Vorsicht würde man nicht darauf rechnen können, dass die oben mit $\beta$ bezeichnete Grösse in der ersten Lage des Fernrohrs bei allen Repetitionen immer denselben, und in der zweiten immer genau den entgegengesetzten Werth behält.

II. Dass die optische Achse des Theodolithenfernrohrs für eine der beiden Ocularstellungen genau berichtigt, d. i. gegen die horizontale Drehungsachse normal sei, ist nicht nöthig für die hier beschriebenen Operationen: dient aber der Theodolith zugleich als Ablesungsfernrohr, so muss allerdings vor solchem Gebrauch diese Berichtigung gemacht sein, und zwar für diejenige Stellung der Ocularröhre, bei welcher beobachtet wird, oder wo Marke und Spiegelbild der Scale deutlich erscheinen. Bekanntlich prüft man die Normalität der optischen Achse zur horizontalen Drehungsachse durch Umlegen, und zwar gerade durch dasjenige Umlegen, welches bei obigen Winkelmessungen \textit{nicht} angewandt werden durfte [S. 439 d. B.], nemlich indem man die Zapfen in die entgegengesetzten Lager legt, ohne den Sinn der Richtung des Fernrohrs zu verändern. Gewöhnlich bezieht sich eine solche Prüfung auf diejenige Stellung der Ocularröhre, wobei man sehr entfernte Gegenstände deutlich sieht. und in diesem Falle ist allerdings weiter nichts nöthig, als dass ein solcher Gegenstand vor und nach dem Umlegen auf dem Fadenkreuze erscheine: in dem gegenwärtigen Falle aber muss man, wenn nach dem Umlegen der vor der Mitte des Objectivs herabhangende Lothfaden eine andere Lage hat als vorher, einen zweiten Zielpunkt neben dem ersten in eben so viel veränderter Lage anwenden. Offenbar muss auch ein an%
\SourceLeaf{449}%
statt des Theodolithen angewandtes besonderes Ablesungsfernrohr derselben Berichtigung unterworfen werden, und also eine dazu taugliche Aufstellung haben; von selbst versteht sich, dass auch die horizontale Drehungsachse gehörig nivellirt sein muss. Die beiden bei den hiesigen Magnetometern gebrauchten Ablesungsfernröhre haben, bei einer bedeutend stärkern optischen Kraft, als man den Theodolithenfernröhren zu geben pflegt, fast ganz dieselbe Aufstellung, wie Theodolithen, nur ohne getheilte Kreise.

Übrigens mag noch bemerkt werden, dass der \textit{Einfluss} eines Fehlers der Collimation auf das Azimuth der optischen Achse von dem Collimationsfehler selbst nur ein sehr kleiner Bruchtheil ist, welcher durch den Unterschied der Secanten der beiden Neigungen bestimmt wird, indem das Fernrohr einmal gegen die Marke, und dann gegen den Spiegel gerichtet ist. Bei dem hiesigen Unifilarmagnetometer sind diese Neigungen $1^0 55'$ und $5^0 16'$: der Unterschied der Azimuthe der optischen Achse, bei der Richtung auf Marke und Spiegel, beträgt folglich nur $\frac{1}{272}$ des Collimationsfehlers selbst. Unter ähnlichen Umständen wird man sich daher gewöhnlich damit begnügen können, die Collimation an einem entfernten Gegenstande zu berichtigen: denn wenn nicht in Folge solcher Berichtigung das Fadenkreuz weit aus der Mitte der Ocularröhre gekommen ist, wird das weitere Herausziehen der letztern schwerlich einen Collimationsfehler erzeugen können, der mehr als einen kleinen Bruchtheil einer Bogenminute betrüge, so dass der Einfluss davon durchaus unmerklich bleibt.

III. Der Zweck, warum man den Lothfaden am Beobachtungsfernrohre fortwährend hängen lässt, besteht darin, dass eine zufällige Verrückung der Scale sofort erkennbar werden soll. Hat eine solche Statt gefunden, so mag man entweder die Scale wieder in ihre vorige Stellung bringen, oder auch in der Rechnung von dem Punkte der Scale, welcher dem Lothfaden nach der Veränderung entspricht, eben so zählen, wie vorher von dem frühern. Bei der gegenwärtig im magnetischen Observatorium angewandten Befestigungsart der Scale an der Säule kommen übrigens zufällige Verschiebungen gar nicht mehr vor.

\clearpage
\fancyhead[L]{\small BEOBACHTUNGEN DER MAGNETISCHEN INCLINATION.}
\SourceLeaf{450}\pdfbookmark[0]{Beobachtungen der magnetischen Inclination in Göttingen.}{inclination-observations-title}
\begin{center}
{\Large\bfseries BEOBACHTUNGEN\par\medskip DER MAGNETISCHEN INCLINATION IN GÖTTINGEN.}
\par\bigskip
Resultate aus den Beobachtungen des magnetischen Vereins. 1841. II.
\end{center}
\Article{1}
Das Inclinatorium, mit welchem die hier mitzutheilenden Beobachtungen angestellt sind, ist von \textsc{Robinson}; es war das letzte Instrument dieser Art, welches der ausgezeichnete Künstler geliefert hat.

Der verticale Kreis hat im Lichten den Durchmesser $241.169$ Millimeter und ist von zehn zu zehn Minuten getheilt; der Abstand zweier Theilstriche an ihren innern Enden beträgt daher $0.351$ Millimeter. Die Theilstriche erscheinen auch im Mikroskop unter beträchtlicher Vergrösserung sehr edel; ihre Breite habe ich durch die an mehrern gemachten Messungen $=0.024$ Millimeter gefunden, so dass einer nahe $41$ Secunden deckt.

Der Durchmesser des horizontalen Kreises, da gemessen, wo die Theilstriche von dem Ende des Indexstriches getroffen werden, ist $148$ Millimeter; die Theilung geht durch halbe Grade und der Vernier gibt einzelne Minuten: es findet nur Eine Ablesung Statt.

Die Grade des Verticalkreises sind von beiden Endpunkten eines horizontalen Durchmessers an nach oben und nach unten bis $90$ gezählt, eine Einrichtung, welche vielleicht in den gewöhnlichen Beobachtungsfällen bequem scheinen mag, aber leicht Verwirrung hervorbringt, wenn man sich einer absichtlich belasteten Nadel bedient, und diese dadurch in einen andern Quadranten tritt, oder wenn man auch Beobachtungen in einer gegen den magnetischen Meridian rechtwinkligen Verticalebene anstellt: wenigstens macht diese Einrichtung in sol%
\SourceLeaf{451}%
chen Fällen eine etwas beschwerlichere und weniger übersichtliche Protocollführung nothwendig. Ich würde daher eine in unverändertem Sinne von $0$ bis $360^0$ oder zweimal von $0$ bis $180^0$ fortlaufende Graduirung vorziehen, und habe mich gewöhnt, immer im untern Quadranten auf der linken, oder im obern auf der rechten Seite anstatt der gravirten Zählung sofort die Ergänzung zu $180^0$ niederzuschreiben: auf diese Art sind in gegenwärtigem Aufsatze alle Ablesungen angegeben. Am horizontalen Kreise laufen die Zahlen zweimal in einerlei Sinn von $0$ bis $180^0$; natürlicher und bequemer wäre eine ununterbrochene Durchzählung bis $360^0$, und in dieser Form habe ich die hier vorkommenden Ablesungen angesetzt.

An der Libelle entspricht ein Ausschlag von einem Millimeter einer Neigung von $9$ Secunden.

\Article{2}
Zu dem Instrumente gehören vier Nadeln, die ich durch die Zahlen $1,2,3,4$ unterscheide: die beiden letzten haben drehbare Achsen, auf welche Einrichtung ich weiter unten zurückkommen werde. An allen acht Zapfen hat die mikroskopische Abmessung keinen Unterschied der Dicke erkennen lassen; ich habe diese Dicke $=0.590$ Millimeter gefunden. Die Nadeln $1$ und $2$ wiegen jede $16.5$ Gramme, die beiden andern jede $20.5$ Gramme.

In den Längen der einzelnen Nadeln finden sich kleine Unterschiede; die Messung ergibt
\begin{center}
\begin{tabular}{crcrl}
für &$1$&$.\quad.\quad.\quad.\quad.\quad.$&$240.931$&Millimeter\\
&$2$&$.\quad.\quad.\quad.\quad.\quad.$&$240.866$&—\\
&$3$&$.\quad.\quad.\quad.\quad.\quad.$&$240.938$&—\\
&$4$&$.\quad.\quad.\quad.\quad.\quad.$&$240.954$&—
\end{tabular}
\end{center}
Die kürzeste der Nadeln ist also nur um $0.303$, und die längste nur um $0.215$ Millimeter kürzer, als der Durchmesser des Kreises im Lichten. Dieser Umstand ist nun zwar dem schärfern Ablesen förderlich, hat aber zugleich die Folge, dass schon eine sehr geringe Excentricität die freie Bewegung der Nadel stören kann, und dass es daher schwer ist, diejenigen Theile des Instruments, von deren Stellungen die Excentricität abhängt, auf eine ganz befriedigende Art zu reguliren, zumal da die Stellungen noch vier andern, zusammen also \textit{sechs} Bedingungen Genüge leisten sollen.
